\clearpage
\thispagestyle{empty}
\begin{center}
    \bfseries\fontsize{14}{16}\selectfont RESUMEN
\end{center}
\vspace{16pt}

\noindent
El \textit{Capacitated Vehicle Routing Problem} (CVRP) es un problema clásico de
optimización combinatoria en el que se debe determinar un conjunto de rutas de costo
mínimo para abastecer clientes con demandas conocidas mediante una flota homogénea de
vehículos con capacidad limitada. En el Informe~Nro.~1 este problema se abordó mediante
métodos exactos de programación matemática, cuyo esfuerzo computacional crece de forma
abrupta con el tamaño de la instancia. El presente informe aborda el CVRP desde la
perspectiva de las \textit{metaheurísticas}, implementando \textit{Particle Swarm
Optimization} (PSO) como método aproximado capaz de obtener soluciones de alta calidad
sin certificado de optimalidad, y contrasta ambos enfoques sobre un conjunto común de
instancias. Para adaptar PSO, un algoritmo de naturaleza continua, al espacio combinatorio
del CVRP, se implementó la representación SR-1 \parencite{ai2009particle}, que decodifica
cada partícula a una lista de prioridad de clientes y puntos de referencia de vehículos,
construyendo rutas factibles mediante inserción de menor costo y mejora local 2-opt. Los
hiperparámetros de PSO se calibraron con Optuna (60 \textit{trials}) sobre instancias de
tuning distintas de las de evaluación. El método se evaluó con 31 repeticiones
independientes sobre nueve instancias de CVRPLIB: las cinco del Informe~1 (para
comparación directa) y cuatro adicionales de mayor tamaño, fuera del alcance del método
exacto. En las cinco instancias compartidas, PSO obtuvo un GAP promedio entre $0{,}03\%$
y $10{,}07\%$ respecto al óptimo, en tiempos de ejecución de fracciones de segundo,
frente a un método exacto que llegó a tardar $4{,}7$ horas sin certificar optimalidad en la
instancia más grande compartida. En las cuatro instancias adicionales, el GAP promedio
se ubicó entre $15{,}25\%$ y $26{,}84\%$, con el máximo en la instancia más grande. La
comparación evidencia una
complementariedad entre ambos enfoques: el método exacto certifica optimalidad mientras
es tratable, pero pierde esa tratabilidad de forma abrupta más allá de cierto tamaño;
PSO no certifica, pero mantiene tiempos acotados y resuelve instancias considerablemente
mayores que las alcanzables por el método exacto.

\vspace{10pt}
\noindent\textbf{Palabras clave:} CVRP; optimización combinatoria; metaheurísticas;
\textit{Particle Swarm Optimization}; PSO; CVRPLIB.
